% ============================================================
% Abstract
% Thesis: Solving the FSRCPSP Using Hybrid ML Approaches
% Author: Souvik Chakraborty
% ============================================================

\begin{center}
  {\Large\bfseries Abstract}
\end{center}
\addcontentsline{toc}{chapter}{Abstract}
\vspace{0.5cm}

Although decomposition heuristics are widely used in project
scheduling, few methods explicitly combine machine learning with
mixed-integer programming to guide how large-scale stochastic
formulations are broken down. In this thesis, we consider a
stochastic resource-constrained project scheduling problem with
flexible resource profiles, where resource allocations may vary
within activity-specific bounds and workloads are modelled as
random variables. The problem is formulated as a mixed-integer
linear programme with chance-constraints that minimises the
project duration while ensuring a predefined confidence level.
Uncertainty in the workloads is approximated through a sample
average approximation based on discrete scenario sets. Exact
solutions are obtainable for small instances, but larger ones
require heuristic approaches. We therefore develop a group-based
activity-oriented Relax-and-Fix heuristic that embeds
unsupervised machine learning into the decomposition strategy.
Rather than fixing one activity per subproblem, the method
extracts a feature vector for each activity, capturing network
topology, scheduling urgency, and stochastic workload statistics,
and applies Ward hierarchical clustering to partition the
activities into groups that are coherent in both structure and
stochastic profile. Each group is then fixed simultaneously in a
single subproblem, giving the solver more combinatorial context
for branching and bounding. A serial schedule generation scheme
provides a warm-start solution, and a stop-on-first-feasible
termination rule allows the procedure to halt as soon as an
integral solution is certified. For instances beyond medium
scale, two further extensions are introduced: time aggregation,
which compresses the time axis by a configurable factor to reduce
variable and constraint counts, and per-group scenario reduction,
which selects a representative subset of scenarios for each
subproblem. All methods are evaluated on standard PSPLIB-based
benchmark instances with respect to solution quality, feasibility
rates, runtime, and the contribution of individual feature groups
to the clustering.

\vspace{0.5em}
\noindent\textbf{Keywords:}
Stochastic project scheduling, Flexible resource profiles, Sample average approximation, Relax-and-Fix, Hierarchical clustering, Mixed-integer programming
